\subsection{Odd terminals with their actual inward path}
\label{sec:odd-complete:inward-anchor}

Write the strictly increasing terminal list as
$(A,b_1,c_1,\ldots,b_j,c_j)$, $A\ge2$, $j\ge1$.
If $b_1\ge4$ use
$F=\bar f_A\prod_i h_{b_i}f_{c_i}$.
If $b_1=3$, the order forces $A=2$; use
$F=\bar f_2f_3h_{c_1}\prod_{i\ge2}h_{b_i}f_{c_i}$.
We prove fourth losses $13/2$ at $pW_a\bar f_b$, $3\le a<b$,
and $6$ at $f_a h_bf_c$, $3\le a<b<c$, keeping these directions
for every positive inner connector $L$.

\paragraph{The first pair has minimum at least four.}
Let $R=r(B)u$ be the one finite sum of the actual pair intervals,
and $\Gamma$ their total scalar loss. Put
$\mu_A=\E\bar f_A$ and
$\sigma_A^2=\operatorname{Var}\bar f_A$.
The actual centered first row of Section~\ref{proof:Q5-C181} is
\begin{equation}\label{eq:odd-anchor:general-first}
\begin{aligned}
 T(\bar f_Ae^R)&\le\bar f_{A+1}+\mu_A BR+\delta,\\
 \delta&=\mu_A\E R+2\sigma_A\|R\|+(12/5)\|R\|^2.
\end{aligned}
\end{equation}
The centered product is bounded in the original measure by
$2\E| (\bar f_A-\mu_A)R|$. The mixed mean and variance use
the same $R$. The three finite-mask records are
\[
\begin{array}{c|c|c|c|c}
 &\Gamma/d&\rho:\|R\|<\rho\sqrt d&\delta_+/d
                   &d^{-1}\sum_{i\ge1}(D_i)_+\\\hline
 j=1,b_1=4&69/20&1/5&72/125&7/25\\
 j=1,b_1\ge5&16/5&3/20&3/5&1/5\\
 j\ge2&16j/5&1/4&5/6&2/5.
\end{array}
\]
Here $D_i=\E B^iR$. The first row uses $A=2,3$ and
$\sigma_A<6\sqrt d/5$; at higher first-pair minimum use
$\sigma_A<3\sqrt d/2$ and its matching stronger mask.

Finite propagation retains the actual inward anchor:
\begin{equation}\label{eq:odd-anchor:general-finite}
 e^\Gamma T^LF\le\bar f_{A+L}+\mu_AB^LR
       +\delta\bar f_{L-1}
       +\mu_A\sum_{i=1}^{L-1}D_i\bar f_{L-1-i}.
\end{equation}
Thus \eqref{eq:odd-complete:input} has $\alpha=-\Gamma$,
$C=1$,
\[
 Q=G_{A+L-1}+\mu_A z^{L-1}r,\qquad
 \E BQu=-\beta_{A+L+1}+\mu_A D_L,
\]
and $c=\delta-\beta_{A+L}$ at $L=1$, or
$c=2\delta_++2\mu_A\sum_{i<L}(D_i)_+-\beta_{A+L}$ at
$L\ge2$. The true anchor mean $\mu_A$ remains in both terms.
Substituting these data into the auxiliary exponent
\eqref{eq:odd-complete:baseline-curved} gives losses greater
than $13/2+10/9$ for one pair and
$15+10/9+64(j-2)/5$ for several pairs. To verify these numbers
directly, the complete scalar and joint-norm bounds are
\[
\begin{array}{c|cc|cc}
 &\multicolumn{2}{c|}{-(\alpha+c+\E BQu-d)/d}
 &\multicolumn{2}{c}{\|(Q-B)u\|/\sqrt d}\\
 &L=1&L\ge2&L=1&L\ge2\\\hline
 j=1,b_1=4&69/20+51/32+1-72/125-1/6&7711/2000&131/80&133/100\\
 j=1,b_1\ge5&16/5+19/16+1-3/5-1/30&889/240&127/80&13/10\\
 j\ge2&16j/5+277/240&16j/5-29/80&27/16&11/8.
\end{array}
\]
There is no auxiliary $\beta_2$ curvature cost; the $10/9$ just
reclaims the isolated uniform fee in the corresponding general
root certificate of Section~\ref{proof:Q5-C181}.

For a minimum-one root these polynomials have nonnegative real
coordinate, imaginary coordinate at most $5/4$ and norm at most
$3\sqrt d/2$. For the coordinate assertions use the actual mask,
whose first degree is at least four:
$|\operatorname{Re}r|,|\operatorname{Im}r/z|
 \le(1024/855)x^2<1/5$.
The real coordinate is at least $\kappa-1/5$, while the
imaginary coordinate is at most $6/5$. The norm follows from
$\|R\|<\sqrt d/4$, $|G_2|<6/5$ and $|G_n|\le1$ for
$n\ge3$, keeping $\mu_A\le1$.
Their $\Phi$ increment is below $3d/5$;
$\Psi(5/4,3/2)<0$. The full-domain loss exceeds $13/2$.
For a positive root write $Q=G_{A+L-1}+\Delta Q$ and use
the integral form \eqref{eq:odd-complete:positive-norm}.
The actual tail norms are at most $\rho\sqrt d$, including
$L=1$ and every finite $A,L$. The increments are at most
$19d/15$, $7d/10$ and $11d/6$ in the three rows.
The weakest resulting loss is
$13/2+10/9-19/15=571/90>6$.

\paragraph{The first pair has minimum three.}
The actual short anchor is $\bar f_2=1+u+Bu$.
Let $c=c_1\ge4$, $B_0=\beta_{c-1}$ and
\[
 r_0=-B_0-z^3G_{c-3}(-z)=-B_0+x^2V+zxU,
 \qquad\Gamma_0=d+\beta_c+B_0\ge35d/16.
\]
At $c=4$, $(U,V)=(1,0)$; thereafter the coordinates of the
same finite $G_{c-3}(-z)$ are used. For one pair take $R=r_0(B)u$,
$\Gamma=\Gamma_0$. For several pairs add the one actual finite
later mask $S$ to $R$ and its true scalar losses to $\Gamma$.
On $d\le1/8$, both cases have the centered row
\begin{equation}\label{eq:odd-anchor:short-first}
\begin{aligned}
 T(\bar f_2e^R)&\le\bar f_3+BR+D,\\
 D&\ge\max\{\E R+\E(uR)+\E((Bu)R)
          +\sqrt{(d+\delta_2)\E R^2}+(12/5)\E R^2,0\}.
\end{aligned}
\end{equation}
The linear coefficient is one, rather than $\mu_2$; all three
mixed means are retained. The actual finite heads $c=4,5$ and
the remaining contracting tail give the following records:
\[
\begin{array}{c|c|c|c|c|c}
 &\Gamma/d&D/d&D_1/d&D_2/d&D_i/d\ (i\ge4)\\\hline
 j=1&35/16&3/5&[-\delta_3/d,d]&[0,\delta_3/d]&\le1/36\\
 j\ge2&431/80+16(j-2)/5&3/2&[-19/60,19/120]&\le1/4&\le1/9.
\end{array}
\]
Also $D_3\le\delta_4$ in the first row and $D_3\le d/6$ in
the second. These are the genuine first-row and signed-tail
certificates of Section~\ref{proof:Q5-C181} for the displayed
short anchor.

The finite propagation is
\begin{equation}\label{eq:odd-anchor:short-finite}
 e^\Gamma T^LF\le Y_L:=\bar f_{L+2}+B^LR+D\bar f_{L-1}
                  +\sum_{i=1}^{L-1}D_i\bar f_{L-1-i}.
\end{equation}
Expanding its first three actual paths gives the exact rows
below for \eqref{eq:odd-complete:input}, with $\alpha=-\Gamma$:
\begin{equation}\label{eq:odd-anchor:short-rows}
\begin{array}{c|c|c|c}
 L&C&c&Q\\\hline
 1&1&D-d&G_2+r\\
 2&1-d+D&D_1-d&G_3+zr\\
 3&1-2d+D+D_1&-\tau+D_2&G_4-d+D+z^2r.
\end{array}
\end{equation}
For example $Y_3$ retains both losses in
$\bar f_5=\one+G_5u-d\bar f_2-dq-\tau\one$; this gives
the displayed coefficient of $q$ and the separate constant
$-\tau+D_2$. Its actual mean is
$-\tau+d^2+D_3-Dd$.

For $j=1$ put $h=2d/5$. The loss coefficients in the same
convolution are $\gamma_3=d-D=h$ and
$\gamma_i=\beta_i-D_{i-3}\ge h$, $i\ge4$.
At $L\ge4$ retain the final four paths $\bar f_0,\ldots,\bar f_3$.
Their sum is $4-d+3u+2Bu+B^2u$, giving
\begin{equation}\label{eq:odd-anchor:single-long}
 C=1-3h,\qquad c=-h(1-d),\qquad
 Q=G_{L+1}+z^{L-1}r-h(2+z),
 \quad\E BQu=-\gamma_{L+3}+2hd.
\end{equation}
For $j\ge2$, the actual first coefficient is instead
$\gamma_3=d-D=-d/2$. Keep its positive contribution
$(D-d)\bar f_{L-1}\le d$, and take $h=d/3$ for every later
loss, including the mean index. Retaining the final three paths
$\bar f_0,\bar f_1,\bar f_2$ gives
\begin{equation}\label{eq:odd-anchor:multiple-long}
 C=1-2h,\qquad c=d-h,\qquad
 Q=G_{L+1}+z^{L-1}r-h,
 \quad\E BQu=-\gamma_{L+3}+hd.
\end{equation}
Both rows follow from \eqref{eq:odd-anchor:short-finite}; their
different constants are kept in the subsequent full vectors.

\paragraph{One complete payment for both root classes.}
The exact curved baseline \eqref{eq:odd-complete:baseline-curved}
and the two transfer lemmas give the following sufficient ledger.
An entry is a loss for $J_0$, followed by an upper bound on its
increment at the stated root. For multiple pairs add
$64(j-2)/5$ to every baseline.
\[
\begin{array}{c|cccc}
 &L=1&L=2&L=3&L\ge4\\\hline
 j=1:\ -J_0/d&1173/100&5573/600&4291/320&799/80\\
 \text{minimum-one increment}&2&2&2&2\\
 \text{positive increment}&2&0&49/12&10/3\\\hline
 j\ge2:\ -J_0/d&12&16&19&16\\
 \text{minimum-one increment}&7/2&7/2&7/2&15/4\\
 \text{positive increment}&29/5&171/50+9/8&25/9+15/8&11/6+3/2.
\end{array}
\]
Here is a direct check of the entries that differ from an
unweighted path input. For $j=1,L=1$, the actual auxiliary joint
vector has squared norm at most
$(1+x^2)^2+x^3<36/25$. The scalar is at most
$-(35/16+91/40)d$, and $C-1=0$. This gives the first loss
$35/4+91/10-153/25=1173/100$.
For the other single-pair rows, the respective scalar upper
bounds are $-1867d/480$, $-211d/40$ and $-403d/80$;
their joint squared-norm caps are $36d/25$, $25d/16$ and
$9d/4$. Keeping the actual curved coefficients from
\eqref{eq:odd-anchor:short-rows}--\eqref{eq:odd-anchor:single-long}
gives the displayed weaker baseline losses. For several pairs
the same calculation uses scalar upper bounds
$-\Gamma-161d/120$, $-\Gamma-329d/480$,
$-\Gamma-159d/80$ and $-\Gamma-31d/24$.

The single-pair polynomials have real coordinate nonnegative,
imaginary coordinate at most $3/2$, norm at most $8\sqrt d/5$,
and $\gamma\ge-8d/5$. Thus
$\Phi(3/2,8/5,1/8)<2$ and
$\Psi(3/2,8/5)<7/10$. The multiple-pair values fit
$\bar D=2$, $R=2$, $\gamma\ge-d$; their respective costs
are below $7/2$ and $15/4$.

For the positive root apply
\eqref{eq:odd-complete:positive-curved} to the same $C,c,Q$.
The single-pair first row has real overshoot at most $x^2$ and
imaginary overshoot at most $x$. Its increment is below $2d$.
The second row has only imaginary overshoot $x^2$, costing
$85d/252$; $\gamma=-2d/5$ contributes less than $3d^2$,
so the saved $d$ covers both on $d\le1/8$.
At the third row, the integral norm increment is at most
$11d/6$, while $|\gamma|\le47d/30$ and
$|\E V_*|\le7d/6$ give a curved correction below $18d^2$.
At longer rows use the same $11d/6$ norm increment,
$\gamma=-6d/5$, $|\E V_*|\le d$ and correction below
$12d^2$. For the latter mean bound, the finite base and root
coordinates differ by at most $x$ each, the actual propagated
tail contributes at most $\max(x/8,x^3)<1/14$, and $h\le1/20$;
$2\theta+1/14+1/20<1$.
For multiple pairs the corresponding norm increments
are $29d/5$, $171d/50$, $25d/9$ and $11d/6$;
the first curved coefficient is zero and the other corrections
are below $9d^2$, $15d^2$ and $12d^2$.
These are comparisons with the complete curved baseline;
$\gamma\one$ and its signed cross term occur in the positive
root's vector. The table leaves more than $13/2$ at a
minimum-one root and more than six at a positive root.

\paragraph{The same short anchor on the high domain.}
For $d\ge1/8$, the actual inequality
\eqref{eq:odd-q:high-three-pair} gives $f_3h_{c_1}\le e^{-d}$.
Each later pair has minimum at least five and is at most
$e^{-2d}$. Hence the same propagated terminal satisfies
\[
 T^LF\le e^{-[1+2(j-1)]d}\bar f_{L+2}.
\]
The bare comparisons of Corollary~\ref{cor:odd-complete:bare}
give rates at least $4+29/10+8(j-1)>13/2$ and
$4+11/2+8(j-1)>6$, respectively.
The closed endpoint $d=1/8$ is shared with the low domain.
At $d=0$ the actual vectors vanish and leaf integration gives one.
